%======================================================================
\subsection{BW-PDM-07 --- Plan: the LE-02001 pilot and the CAD register (24 Oct 2026, approved 24 Oct)}\label{sec:pdm07-plan}
%======================================================================

\eref{PDM-07} is the first real release: the Leonardo pedestal \code{LE-02001} and its parts go through the vault, the
review and the release, and come out as rev~A with drawings and a release package, with no Fusion step. It is the exit
test of M2, so it exercises almost every epic of the PDM, drawing and assembly tracks. This plan sets out the pilot as a
sequence with a checklist, what each step exercises, what is still missing, and the facts found about the Leonardo files
and register. The acceptance stays as written: \emph{\code{LE-02001} rev~A released entirely from FreeCAD, with drawings
and package.} It depends on \eref{PDM-05} (building) and \eref{DRW-05} (not started).

%----------------------------------------------------------------------
\subsubsection*{1. What was found (24 Oct)}

A read-only search of this Mac found two Leonardo repositories in \file{\textasciitilde/Documents/GitHub/}, and no
FreeCAD or Fusion model of the pedestal.
\begin{center}
\small
\begin{tabularx}{\linewidth}{@{}>{\raggedright\arraybackslash}p{3.6cm}X@{}}
\toprule
Where & What it holds\\
\midrule
\file{Leonardo}, branch \code{main} & READMEs only under \file{cad/source/LE-02001},
  \file{cad/exports}, \file{cad/register}, \file{drawings}: ``Live design: Fusion project \emph{Leonardo} / folder
  \emph{v0\_proto}''. No \file{.f3d}, \file{.f3z}, STEP or PDF of an LE part.\\
\file{Leonardo-restructure}, restructure branch &
  Branch \code{restructure/docs-designs-references}. ADR-0002 (proposed 16 Oct, awaiting you): \file{designs/register/register.csv} becomes the sequence log, checked by
  \file{tools/numbering/check\_ids.py}; the Notion \emph{CAD \& Drawing Register} becomes its mirror; scripted parts
  (build123d, CadQuery) have the script as master. STEP files there are scripted parts (link plates, cycloid, housing),
  not the pedestal.\\
\file{register.csv} (restructure) & 12 issued rows, all at \code{A.1}, source \code{fusion}: \code{LE-02001} Pedestal Assy and 11 parts
  whose parent is \code{LE-02001}. 6 blocked rows, \code{LE-01012}--\code{LE-01017} (test-arm placeholder names).
  No COTS rows yet.\\
FreeCAD test arm & \file{tools/freecad/test\_arm/leonardo\_test\_arm.FCStd}: the phase-0 test arm, with placeholder names.\\
LE-STD-0001 Rev A.2 & The numbering standard (below). Rev A.1 is on \code{main}, Rev A.2 on the restructure branch.\\
\bottomrule
\end{tabularx}
\end{center}
So the pedestal exists only in Fusion, and the first question is where the pilot's files come from.

\paragraph{The numbering standard and the vault disagree in four places.} The spike of \S\ref{sec:pdm06-plan} also
parsed LE numbers with the vault's numbering code:
\begin{center}
\small
\begin{tabularx}{\linewidth}{@{}lXX@{}}
\toprule
 & LE-STD-0001 Rev A.2 & The vault today\\
\midrule
Part number & \code{\{PROG\}-\{BB\}\{SSS\}}, band + sequence: \code{LE-01001} & \code{\{PROG\}-\{DT\}-\{PPPPP\}-\{SSS\}}:
  \code{LE-01-01001-001}. \code{LE-01001-A.1} does not parse, so an import gives it a new number; the first assembly
  becomes \code{LE-02-01001-001}, not \code{02001} (spike).\\
Working revisions & \code{A.1}, \code{A.2}; after release, \code{B.1}, \code{B.2} towards B & \code{01}, \code{02};
  after release, \code{B} in WIP under an ECO (your answer, 23 Oct)\\
Skipped letters & I, O, Q & I, O, Q, S, X, Z (ASME Y14.35); \code{S}, \code{X}, \code{Z} are refused (spike)\\
Drawing number & the part's number (``part number = drawing number'') & a drawing is its own item with its own number
  (open question \ref{q:dwgpn})\\
\bottomrule
\end{tabularx}
\end{center}
LE-STD-0001 also has bands the vault's default document types lack (\code{03} tooling, \code{04} harness, \code{05}
electronic, \code{08} reference, \code{09} COTS) and a register class per band. None of this blocks the pilot, but each
needs a decision before the first LE item is created, because numbers never change afterwards.

\begin{figure}[!htbp]
\centering
\begin{tikzpicture}[x=1cm,y=-0.47cm,
  it/.style={font=\scriptsize,anchor=west,inner sep=1pt},
  pbs/.style={font=\tiny\ttfamily\color{bwgrey},anchor=west}]
\node[it,font=\scriptsize\bfseries] (top) at (0,0) {LE-02001\quad Pedestal Assy};
\foreach \i/\pn/\t/\p in {1/LE-01001/Base plate/PED, 2/LE-01002/Column \slash{} riser/PED,
  3/LE-01003/J1 yaw joint module/ACT, 4/LE-01004/Top deck \slash{} J2 interface/PED, 5/LE-01005/Rear EE shelf/PED,
  6/LE-01006/E-stop mushroom bracket/SAF, 7/LE-01007/Cable guides \slash{} strain relief/HAR,
  8/LE-01008/J2 pitch joint module/ACT, 9/LE-01009/Intermediate EE plate/PED, 10/LE-01010/Lower vertical supports/PED,
  11/LE-01011/Upper vertical supports/PED} {
  \draw[black!35] (0.15,\i-0.6) |- (0.5,\i);
  \node[it] at (0.55,\i) {\pn\quad \t};
  \node[pbs] at (6.2,\i) {\p};
}
\node[pbs] at (6.2,0) {PBS};
\node[font=\tiny\color{bwgrey},align=left,anchor=north west] at (7.3,0.4) {All rows: rev \texttt{A.1}, status issued,\\
source \texttt{fusion}.\\[3pt] \texttt{LE-01003}, \texttt{LE-01008}: class part,\\``holds the drive until it is split into\\an
02xxx assembly''.\\[3pt] \texttt{LE-01005}: ``demote pending;\\EE moved to \texttt{LE-01009}''.\\[3pt]
Blocked, never issued:\\\texttt{LE-01012}--\texttt{LE-01017}.};
\end{tikzpicture}
\caption{\code{LE-02001} as the register has it today (\file{register.csv}, restructure branch): one level, eleven
parts, no hardware. The roadmap's stack order \code{LE-01001} $\rightarrow$ \code{01010} $\rightarrow$ \code{01009}
$\rightarrow$ \code{01011} $\rightarrow$ \code{01002} is the order of assembly (joints), not sub-assemblies.}
\label{fig:pdm07-register}
\end{figure}

%----------------------------------------------------------------------
\subsubsection*{2. The pilot, step by step}

\begin{figure}[!htbp]
\centering
\begin{tikzpicture}[>=Stealth,x=1cm,y=-1.02cm,
  st/.style={draw=black!30,fill=white,rounded corners=2pt,font=\scriptsize,anchor=west,minimum height=0.72cm,text width=6.1cm,align=left},
  n/.style={circle,fill=bwdark,text=white,font=\tiny\bfseries,inner sep=0pt,minimum size=11pt}]
\foreach \y/\txt in {
  0/{Decide: files, number form, revisions, register master (Q1--Q5)},
  1/{LE vault in iCloud Drive: program LE, bands, standards, properties},
  2/{Register in: 12 issued rows become items, 6 blocked numbers reserved},
  3/{Part files in: one file per part, properties filled, checked in},
  4/{\code{LE-02001}: joints, BOM against a hand count, where-used},
  5/{Drawings: 11 parts + the assembly, parts list, revision block},
  6/{Design freeze (G1): checklist and validation macro, report kept},
  7/{Review and release, children first, then \code{LE-02001}: rev A},
  8/{Publish PDFs, build the release package with checksums},
  9/{Register and Ops: rows to released A; BOM to the outbox},
  10/{Afterwards: ECR $\rightarrow$ ECO $\rightarrow$ rev B on \code{LE-01001}}} {
  \node[n] at (0,\y) {\y};
  \node[st] at (0.3,\y) {\txt};
}
\def\chips#1#2{\node[anchor=west,inner sep=0pt] at (6.95,#1) {#2};}
\def\ck#1#2{{\setlength{\fboxsep}{1.6pt}\colorbox{#1}{\tiny\bfseries\ttfamily\textcolor{white}{#2}}}}
\chips{0}{\ck{bwgrey}{you}}
\chips{1}{\ck{bwgreen}{PDM-01}\,\ck{bworange}{PDM-03}\,\ck{bworange}{DRW-01}}
\chips{2}{\ck{bwgrey}{new: register link}}
\chips{3}{\ck{bworange}{PDM-04}\,\ck{bwgrey}{ASM-01}\,\ck{bworange}{ASM-04 props}}
\chips{4}{\ck{bworange}{ASM-03}\,\ck{bworange}{ASM-04}\,\ck{bwgrey}{PDM-06}}
\chips{5}{\ck{bworange}{DRW-01}\,\ck{bwgrey}{DRW-02}\,\ck{bwgrey}{DRW-05}}
\chips{6}{\ck{bwgrey}{new: freeze}\,\ck{bwgrey}{ASM-06}}
\chips{7}{\ck{bworange}{PDM-05}}
\chips{8}{\ck{bwgrey}{DRW-08}\,\ck{bwgrey}{MFG-04}}
\chips{9}{\ck{bwgrey}{new: register sync}\,\ck{bwgrey}{PDM-06}}
\chips{10}{\ck{bworange}{PDM-05}}
\node[font=\tiny,anchor=west] at (0.3,11.0) {\ck{bwgreen}{\phantom{x}} built\quad \ck{bworange}{\phantom{x}} in progress\quad
  \ck{bwgrey}{\phantom{x}} not started (or new in this plan)};
\end{tikzpicture}
\caption{The pilot as a sequence. Each step names the epics it exercises, coloured by their state on 24 Oct. Step~7 needs
every step before it; steps 5 and 8 hold the most that is not built.}
\label{fig:pdm07-flow}
\end{figure}

\providecommand{\tick}{\tikz[baseline=-0.6ex]\draw[black!50] (0,-0.9ex) rectangle (1.8ex,0.9ex);}
\begin{center}
\small
\begin{longtable}{@{}p{0.4cm}>{\raggedright\arraybackslash}p{5.6cm}>{\raggedright\arraybackslash}p{6.8cm}p{0.5cm}@{}}
\toprule
\# & Do & Check (and what it proves) & OK\\
\midrule\endhead
\bottomrule\endfoot
1 & Create the LE vault in \file{iCloud Drive/BruceWorx Vaults/} with program \code{LE} and the bands as document
  types; drawing standards; properties \code{PBS}, \code{Source}, \code{Vendor}, \code{VendorPartNumber} &
  The Create Vault panel does it without the command line; \code{LE-01001} is accepted as a number (Q2) & \tick\\
2 & Read \file{register.csv}: create the 12 issued items, mark the 6 blocked numbers used & Next free part is
  \code{LE-01018}, as \code{check\_ids.py --next 01} says; nothing is renumbered & \tick\\
3 & Put each part's file in (model, or import its STEP, Q1); fill Description and material; check in with a comment &
  Vault panel shows 11 parts, WIP, checked in; mass is computed for every part & \tick\\
4 & Build \code{LE-02001} from the parts in stack order with joints; add the fasteners as COTS items (Q9) &
  Opens from a moved vault with no broken links (\eref{ASM-01}); BOM equals a hand count; where-used of a screw lists
  \code{LE-02001} (\eref{PDM-06}) & \tick\\
5 & One drawing per part and one for the assembly, from the vault's template & Title blocks fill themselves
  (\eref{DRW-01}); parts list matches the BOM (\eref{DRW-02}); revision block and PRELIMINARY stamp (\eref{DRW-05}) & \tick\\
6 & Run \emph{Design freeze\ldots} on \code{LE-02001} & Every automatic line passes or is waived with a reason; the
  report is kept with the release (section~3) & \tick\\
7 & Submit and approve the parts and drawings, then the assembly (Q11) & Each approval records its checklist; files turn
  read-only; Browser shows \code{A} & \tick\\
8 & Publish and package & PDFs print \code{LE-01001-A}; the package (STEP AP242, PDF, 3MF) rebuilds byte-identically
  (\eref{MFG-04}) & \tick\\
9 & Sync the register & \file{register.csv} rows read \code{A}, released; the diff is yours to commit; the Ops outbox
  holds the BOM & \tick\\
10 & Raise an ECR against \code{LE-01001}, open the ECO, change it, release B & \eref{PDM-05}'s acceptance on a real
  part; \code{LE-02001} A still lists \code{LE-01001} A (pinned, \S\ref{sec:pdm06-plan}) & \tick\\
\end{longtable}
\end{center}

%----------------------------------------------------------------------
\subsubsection*{3. Theory: the release closure and the order of release}

To release $t = $ \code{LE-02001} the vault must release its \emph{closure}: every item reachable from $t$ in the BOM
graph of \S\ref{sec:pdm06-plan}, plus each item's drawing,
\[
  C(t) = D(t) \cup \{\,\mathrm{dwg}(x) : x \in D(t)\,\},\qquad D(t) = \{t\} \cup \{\,x : t \rightsquigarrow x\,\}.
\]
The release checklist allows an assembly to be approved only when its components are released (\eref{PDM-05}, line~3),
so the order of approvals is a \emph{reverse topological order} of the BOM graph: leaves first. Grouping items by their
height $h(x)$ (0 for a leaf, $1 + \max h(\text{child})$ otherwise) gives \emph{waves} that can be approved in any
order inside a wave:
\[
  \text{wave}_k = \{\,x \in D(t) : h(x) = k\,\},\qquad k = 0, \ldots, h(t).
\]
For the register's tree (figure~\ref{fig:pdm07-register}), $h(t)=1$: wave~0 is the 11 parts and the fasteners, wave~1 is
\code{LE-02001}. With drawings that is $|C(t)| = 2\cdot12 = 24$ items, plus the COTS items (which have no drawing), and
each needs a submit and an approval: about fifty actions by hand. A \emph{Release tree\ldots} command that walks the waves
and shows one checklist after another keeps that to one pass (question~11). The same walk gives a readiness figure for
the Start page, $\lvert\{x \in C(t) : \text{ready}(x)\}\rvert / \lvert C(t)\rvert$.

\paragraph{The design-freeze checklist (G1).} Your Notion gate ``LE G1 --- Spec lock / CAD freeze ready'' holds the
criteria, which are not readable from here (question~7). Proposed, split by who checks:
\begin{center}
\small
\begin{tabularx}{\linewidth}{@{}lXl@{}}
\toprule
Line & Check & By\\
\midrule
Recompute & Every document of $D(t)$ recomputes with no invalid or touched object & macro\\
Links & Every link resolves inside the vault; no absolute path; no file outside the vault & macro\\
BOM & Every line is a vault item; no cycle; COTS lines have vendor and MPN & macro\\
Properties & Every property the schema marks \code{required} for release has a value (\file{partprops.py}) & macro\\
Register & Every item of $D(t)$ has an issued register row whose title matches its description & macro\\
Interference & No clash between occurrences (\eref{ASM-06}), or each one waived with a reason & macro, when built\\
Drawings & Each item of $D(t)$ that needs one has a drawing whose title block shows its P/N (\file{release\_check.py}) & macro\\
Mass & Total mass $\sum_x Q_t(x)\,m_x$ against the LE-PAR-0001 budget & you\\
Spec lock & The SPC and ICD revisions the design answers to are issued & you\\
\bottomrule
\end{tabularx}
\end{center}
The macro runs headless (\code{FreeCADCmd}) on the checked-in versions, so its result does not depend on what happens to
be open; its report (JSON, and a one-page PDF) is attached to the release as the approval's attachment, which the
journal already supports.

%----------------------------------------------------------------------
\subsubsection*{4. Gaps: what rev A needs that is not built}

\begin{center}
\small
\begin{tabularx}{\linewidth}{@{}llX@{}}
\toprule
Needed & Epic, state & Without it\\
\midrule
Revision block, PRELIMINARY stamp & \eref{DRW-05}, not started & drawings print no revision history: not a real release\\
Parts list and balloons & \eref{DRW-02}, not started & assembly drawing without a parts list; BOM attached as CSV\\
Publish PDF & \eref{DRW-08}, not started & TechDraw's own export, by hand, sheet by sheet\\
Release package & \eref{MFG-04}, not started & a zip made by hand: fails ``entirely from FreeCAD''\\
BOM per revision, where-used & \eref{PDM-06}, not started (\S\ref{sec:pdm06-plan}) & hand count only\\
Indented and flattened BOM, CSV & \eref{ASM-04}, in progress & the Assembly BOM spreadsheet, with its merge bug\\
One part per file tooling & \eref{ASM-01}, not started & the files must already be one per part\\
LE number form, register & \eref{PDM-03}, in progress; register link new & numbers typed by hand (Q2, Q5)\\
Interference & \eref{ASM-06}, not started & that freeze line is skipped\\
Hardware library & \eref{MOD-05}, not started & fasteners modelled and numbered by hand\\
\bottomrule
\end{tabularx}
\end{center}

%----------------------------------------------------------------------
\subsubsection*{5. Build order and effort}

The pilot mostly \emph{uses} other epics; what it builds itself is small.
\begin{center}
\small
\begin{tabularx}{\linewidth}{@{}lXl@{}}
\toprule
\# & Step & Effort\\
\midrule
1 & LE-STD-0001 numbers in the vault (per-program format, bands as types, revision display; with \eref{PDM-03}) & 1.5\,d\\
2 & Register link: read \file{register.csv} (issued, blocked, next free), create items, write state and revision back
  as a diff & 1.5\,d\\
3 & LE vault set-up, part files in (Q1), \code{LE-02001} and its BOM checked by hand & 1\,d\\
4 & Twelve drawings from the template (with \eref{DRW-01}) & 2\,d\\
5 & \emph{Design freeze\ldots}: the headless validation macro, its report, the manual lines & 2\,d\\
6 & \emph{Release tree\ldots}: waves, one checklist per item (optional, Q11) & 1\,d\\
7 & Dry run on a copy of the vault; the real release; register sync; write-up & 1\,d\\
\midrule
 & Total (excluding the epics in section~4 and any modelling) & 10\,d\\
\bottomrule
\end{tabularx}
\end{center}

\paragraph{Risks.} \emph{The files.} If the pilot imports STEP from Fusion, the parts are dumb solids: a change during
the pilot means direct editing (\eref{MOD-04}) or remodelling. \emph{Number decisions are permanent}: once an LE item is
created in the vault its number never changes, so Q2--Q4 must be settled first; the dry run on a copy is there to catch
anything missed. \emph{Two registers} (the CSV and Notion) can disagree; ADR-0002 is still proposed. \emph{Dependencies}:
four of the epics in section~4 are not started; the real release waits for them, the dry run does not.

%----------------------------------------------------------------------
\paragraph{Your answers to the questions below (24 Oct).} Q1 changed: \textbf{start fresh, no importing of files}
(``for H, lets start fresh, no importing files''). \code{LE-02001} and its parts are \textbf{modelled from scratch in
FreeCAD}; nothing is exported from Fusion as STEP, so the dumb-solid risk above falls away and the modelling becomes
part of the pilot (it was left out of the 10\,d). That modelling needs \textbf{your design data} (dimensions, materials,
hardware, the Fusion model to read from): no part is invented, and every number follows LE-STD-0001. Q2--Q11 as
recommended. 2: a \textbf{per-program number format}, LE keeps \code{LE-01001}. 3: the vault's lifecycle with a
per-program revision \emph{display} (\code{A.}$n$, \code{B.}$k$); LE-STD-0001 is amended to skip S, X and Z too. 4: a
drawing stays its own item but carries and prints its part's number. 5: \textbf{ADR-0002 accepted}: the CSV is the
sequence log, the vault writes state and revision back as a diff you commit, Notion mirrors the CSV. 6: each register
row shows \code{P/N-REV}, state and \emph{Open in FreeCAD}. 7: the proposed design-freeze lines, plus anything G1 has
that they miss. 8: a \textbf{dry run now} on a copy of the vault; the real release once \eref{DRW-05}, \eref{DRW-08}
and \eref{MFG-04} are built. 9: screws, inserts and bearings numbered as COTS items (\code{09xxx}). 10: \code{LE-01003}
and \code{LE-01008} released as parts now, split later through an ECO. 11: \emph{Release tree\ldots} is added.

\paragraph{Progress (24 Oct, evening): the Leonardo actuator in your vault.} Ahead of the pedestal, the Leonardo
actuator --- a NEMA~17 motor with a planetary gearbox --- is being designed in your real vault \emph{Vector}, program
LE, by the lead agent driving your FreeCAD through the Claude Code bridge (\S\ref{sec:aut03-plan}): the vault's first
real use. All items are WIP~01: a \emph{Gearbox} file holding the \emph{Parameters} sheet (\S\ref{sec:mod15-plan}) and
the top assembly \code{LE-02-01001-001}; the motor adapter plate \code{LE-01-01001-001}, end cap \code{-01002},
housing with the internal ring (48 teeth) \code{-01003}, planet carrier \code{-01004}, planet (18 teeth)
\code{-01005} and sun (12 teeth) \code{-01006}. Two design decisions were taken with you:
\begin{itemize}[nosep]
  \item \textbf{Pressure angle $25^\circ$}, not $20^\circ$. A rack-cut gear with fewer than
        $z_{\min} = 2/\sin^2\alpha$ teeth is undercut: $z_{\min} = 17.1$ at $\alpha = 20^\circ$, so the 12-tooth sun
        would be undercut, but $z_{\min} = 11.2$ at $25^\circ$, so 12 teeth need no profile shift.
  \item \textbf{Housing, plate and end cap $\varnothing 66$} with through-bolts on a $\varnothing 58$ circle. The smaller
        option, $\varnothing 60$ with bolts on $\varnothing 52$, would put the bolt holes into the ring's teeth (root
        circle $\varnothing 50.3$).
\end{itemize}
The numbers are in the vault's present form (\code{LE-01-01001-001}), not the per-program \code{LE-01001} form of
answer~2, which is not built yet. Gaps met on the way are listed under \eref{FND-08}.

\subsubsection*{Questions before PDM-07 starts}

\begin{enumerate}
  \item \textbf{Where are the \code{LE-02001} and \code{LE-01001} files?} Neither Leonardo repository holds a model of
        the pedestal (section~1); the live design is in Fusion (\emph{Leonardo / v0\_proto}). Is there a FreeCAD model
        anywhere (M1 asks for one modelled from scratch), or should the pilot run on STEP exports from Fusion?\\
        \emph{Recommendation:} tell us the path if FreeCAD files exist. If not, export the 12 items from Fusion as STEP,
        one file per component, into a staging folder, and run the pilot on those; remodelling natively stays M1's exit,
        not the pilot's.
  \item \textbf{Part-number form for LE.} Keep LE-STD-0001's \code{LE-01001} or map to the vault's
        \code{LE-01-01001-001}?\\
        \emph{Recommendation:} the vault learns a per-program format, so LE keeps \code{LE-01001} (band = document type,
        sequence 001--999) and nothing the team already uses changes. Other programs keep the current form.
  \item \textbf{Working revisions.} LE-STD-0001 writes \code{A.1}, \code{A.2} and \code{B.1}; the vault uses
        \code{01}, \code{02} and \code{B} in WIP (your 23 Oct answer). The skipped letters also differ.\\
        \emph{Recommendation:} keep the vault's lifecycle and add a per-program revision \emph{display}: for LE, pre-release
        \code{0}$n$ prints as \code{A.}$n$ and the $k$-th review round of an ECO towards B prints as \code{B.}$k$. The
        mapping is one to one, so nothing is lost. Amend LE-STD-0001 to skip S, X and Z as well (no LE item has passed A).
  \item \textbf{Drawing numbers.} LE-STD-0001 gives a drawing its part's number; the vault makes a drawing its own item
        (open question \ref{q:dwgpn}).\\
        \emph{Recommendation:} a drawing stays its own item (its file and revision are tracked) but carries and prints its
        part's number; internally its key is the part's number with a drawing suffix.
  \item \textbf{Which register is the master.} \file{register.csv} (ADR-0002, proposed), Notion, or the vault?\\
        \emph{Recommendation:} accept ADR-0002. The CSV stays the sequence log: the vault reads it to allocate numbers and
        never issues a blocked one. For items it holds, the vault writes state and revision back as a diff you commit;
        Notion mirrors the CSV.
  \item \textbf{``CAD register links in the BruceWorx app''.} What should a link do?\\
        \emph{Recommendation:} each register row shows the vault's \code{P/N-REV} and state, and an \emph{Open in
        FreeCAD} link that opens the item from the vault; the Ops register tile counts the same rows once Ops reads them.
  \item \textbf{The design-freeze checklist.} The G1 criteria are in Notion. Use the proposed lines (section~3), or send
        the gate's list?\\
        \emph{Recommendation:} the proposed lines, plus anything G1 has that they miss.
  \item \textbf{Dry run first.} Run the pilot now on a copy of the vault, gaps and all, or wait for \eref{DRW-05},
        \eref{DRW-08} and \eref{MFG-04}?\\
        \emph{Recommendation:} dry run now (it finds the next gaps early and releases nothing for real); the real
        release when those three are built.
  \item \textbf{Hardware.} Number the screws, inserts and bearings of \code{LE-02001} as COTS items (\code{09xxx})?\\
        \emph{Recommendation:} yes; the BOM hand count and \eref{PDM-06}'s acceptance need them.
  \item \textbf{Joint modules.} \code{LE-01003} and \code{LE-01008} hold drives but are parts in the register. Split them
        into \code{02xxx} assemblies before the pilot?\\
        \emph{Recommendation:} no; release them as parts now and split them later through an ECO.
  \item \textbf{Releasing a tree.} Approve the 24 parts and drawings one by one, or add \emph{Release tree\ldots}
        (children first, one checklist per item, one pass)?\\
        \emph{Recommendation:} add it; it is a day's work and every later release uses it.
\end{enumerate}
